\documentclass[10pt,a4paper]{amsart}
\usepackage[margin=19mm]{geometry}
\usepackage{amsmath,amssymb,mathtools}
\usepackage[scr=rsfs,cal=euler]{mathalfa}
\usepackage{xfrac}
\usepackage[unicode,hidelinks,bookmarks=false]{hyperref}
\newcommand{\ie}{i.e.\ }
\newcommand{\cf}{cf.\ }
\newcommand{\resp}{respectively\ }
\newcommand{\Subsection}{\subsection}
\providecommand{\nobreakdash}{\nobreak\hbox{-}}
\setlength{\emergencystretch}{2em}
\multlinegap0pt
\usepackage{bm}
\usepackage{bbm}
\usepackage{dsfont}
\usepackage{xspace}
\usepackage{cancel}
\usepackage{mathtools}
\usepackage{amsthm}
\makeatletter
\@ifundefined{theorem}{\newtheorem{theorem}{Theorem}}{}
\@ifundefined{lemma}{\newtheorem{lemma}[theorem]{Lemma}}{}
\makeatother
\newcommand{\ess}{\mathrm{ess}\,}
\title[Spectral theoretic characterisation of Markov chain converge]{Spectral theoretic characterisation of Markov chain convergence}
\author{Bryn Davies, Angelica Yu Xiao}
\date{Source: arXiv:2410.15829v2 (CC BY 4.0)}
\begin{document}
\maketitle

\noindent\emph{Source and licence.} Spectral theoretic characterisation of Markov chain convergence. Version arXiv:2410.15829v2 (CC BY 4.0), distributed under the Creative Commons Attribution 4.0 International licence, as recorded in the arXiv licence metadata for this exact version and shown on the abstract page https://arxiv.org/abs/2410.15829v2. Full rights notice: licenses/arxiv-2410.15829-CC-BY-4.0.txt.\par\smallskip

\noindent\textit{Editorial scope.} Counted unit: the Lemma whose source label is \texttt{lm:Rstep\_approx} in the article (source0.tex lines 464--470, source label \texttt{lm:Rstep\_approx}), delivered together with the article's own complete original proof of it (source0.tex lines 471--511, one \texttt{proof} environment of 41 lines). No mathematics is written, repaired or re-derived: statement and proof are the article's own text line for line. Required context retained uncounted: none. Every definition, display and label of those lines is retained unchanged. Omitted from this package and not redistributed: 1--463, 512--1074. Nothing inside a retained excerpt is omitted or altered: every retained body is delivered exactly as the article's lines are written, so no omission receipt of any kind accompanies this package. Citations: the retained excerpts carry no citation command at all, so no bibliography entry of the article is transcribed and no reference is redistributed. Reference adaptation: none was needed and none was made; every reference inside the delivered unit resolves inside it, so no unresolved reference or citation is produced. Printed slips kept literal: no printed word, symbol, hypothesis or number of the counted statement or of its proof was altered. Layout and class: the article's own class is \texttt{article} with packages including babel, inputenc, fontenc, graphicx, caption, subcaption, appendix, tikz, pgfplots, cutwin, enumerate, geometry, amsmath, yhmath; the wrapper is the lane's pinned \texttt{amsart} editorial wrapper at 10pt on A4, which restates the theorem-like environments the retained text needs and adds the article's own macro definitions that the retained text needs (\texttt{\textbackslash ess}), with extra wrapper packages (bm, bbm, dsfont, xspace, cancel, mathtools). The delivered numbering is the unit's own. Assets: no retained span contains a figure environment, a graphics-inclusion command, a PDF-page inclusion, a table or a TikZ picture; no image file is copied into this package, the package declares no asset dependency and no image file is redistributed with it. Licence: version arXiv:2410.15829v2 of this article is distributed under the Creative Commons Attribution 4.0 International licence, as recorded in the arXiv licence metadata for this exact version and shown on the abstract page https://arxiv.org/abs/2410.15829v2. A full rights notice is delivered at licenses/arxiv-2410.15829-CC-BY-4.0.txt. Characters: every retained excerpt is pure ASCII, so the delivered engine, XeLaTeX with its default Latin Modern text font, reports no missing character.
\begin{lemma} \label{lm:Rstep_approx}
    For every $p\in BV[0,\infty)\cap L^1[0,\infty)$ and $l\in\mathbb{N}$, there exists $p_l$ a rational step function of width $\frac{1}{l}$ on $[0,\infty)$ such that
    \begin{equation*}
        \lVert p_l - p\rVert_{ L^1[0,\infty)} \leq \frac{\ess V(p)}{l},
    \end{equation*}
    where $\ess V(p)=\inf_{\Tilde{p}=p\text{ a.e.}}V(\Tilde{p})$ is the essential variation of $p$ on $[0,\infty)$.
\end{lemma}

\begin{proof}
    Without loss of generality, we write $p$ as the representative of $[p]$ with the smallest variation $V(p) = \ess V([p])$. Let $V(p)_x$ be its variation over $[0,x]$:
    \begin{equation*}
        V(p)_x = \sup_{\mathcal{P}_x}\sum_{i=0}^{N_{\mathcal{P}_x}} |p(a_{i+1})-p(a_i)|,
    \end{equation*}
    where $\mathcal{P}_x$ is a partition of $[0,x]$. Then $V(p)_\infty=V(p)$.
    
    Let $p_1(x) = \frac{1}{2}(V(p)_x+p(x))$ and $p_2(x) = \frac{1}{2}(V(p)_x-p(x))$. Then $p_1$ and $p_2$ are increasing, and $p=p_1-p_2$. Moreover,
    \begin{equation*}
        \lim_{x\rightarrow\infty}p_1(x) = \lim_{x\rightarrow\infty}p_1(x) = \frac{1}{2}V(p).
    \end{equation*}
    This follows from the fact that $\lim_{x\rightarrow\infty}p(x) = 0$. Since $p\in BV[0,\infty)$, $p(x)$ must converge as $x\rightarrow\infty$ (the number of upcrossings of $p(x)$ between any $a<b$ must be finite). Since $p\in L^1[0,\infty)$, the limit must be zero.

    Fix some $l\in\mathbb{N}$. We can approximate the increasing functions $p_1$ and $p_2$ respectively by rational step functions $\Tilde{p_1}$ and $\Tilde{p_2}$ of width $\frac{1}{l}$.
    \begin{align*}
        \Tilde{p_1}(x) &= \sum_{i=0}^\infty p_1(x_i)\cdot\mathds{1}_{[\frac{i}{l},\frac{i+1}{l})},\\
        \Tilde{p_2}(x) &= \sum_{i=0}^\infty p_2(x_i)\cdot\mathds{1}_{[\frac{i}{l},\frac{i+1}{l})},
    \end{align*}
    where each $x_i$ is an arbitrary point in $[\frac{i}{l},\frac{i+1}{l})$.

    Since $p_1$ is increasing, it holds that $|p_1(x)-p_1(x_i)|\leq p_1(\frac{i+1}{l})-p_1(\frac{i}{l})$ for all $i\in\mathbb{N}$ and $x\in[\frac{i}{l},\frac{i+1}{l})$. Therefore,
    \begin{equation*}
        \int_{[\frac{i}{l},\frac{i+1}{l})}|p_1(x)-p_1(x_i)|dx\leq \frac{1}{l}\left(p_1\left(\frac{i+1}{l}\right)-p_1\left(\frac{i}{l}\right)\right)\quad\text{for all $i\in\mathbb{N}$.}
    \end{equation*}
    The sum over $i\in\mathbb{N}$ is telescoping:
    \begin{align*}
        \lVert p_1-\Tilde{p_1}\rVert_{ L^1[0,\infty)} &= \sum_{i=0}^\infty\int_{[\frac{i}{l},\frac{i+1}{l})}|p_1(x)-p_1(x_i)|dx\\
        &\leq\sum_{i=0}^\infty\frac{1}{l}\left(p_1\left(\frac{i+1}{l}\right)-p_1\left(\frac{i}{l}\right)\right)\\
        &=\lim_{i\rightarrow\infty}\frac{1}{l}(p_1(i)-p_1(0))\\
        &= \frac{1}{l}(\frac{1}{2}V(p)-p(0)).
    \end{align*}
    The same result applies for $p_2$:
    \begin{equation*}
        \lVert p_2-\Tilde{p_2}\rVert_{ L^1[0,\infty)}= \frac{1}{l}(\frac{1}{2}V(p)+p(0)).
    \end{equation*}

    Let $p_l = \Tilde{p_1} - \Tilde{p_2}$. Then $p_l$ is also a rational step function of width $\frac{1}{l}$ and we have the bound
    \begin{equation*}
        \lVert p-p_l\rVert_{ L^1[0,\infty)} \leq \lVert p_1-\Tilde{p_1}\rVert_{ L^1[0,\infty)} + \lVert p_2-\Tilde{p_2}\rVert_{ L^1[0,\infty)} \leq \frac{V(p)}{l}.
    \end{equation*}
\end{proof}

\end{document}
